\chapter*{INTRODUCCIÓN}
\addcontentsline{toc}{chapter}{INTRODUCCIÓN}

En el Perú, la máxima demanda del Sistema Eléctrico Interconectado Nacional aumentó de \SI{4579}{\mega\watt} en 2010 a \SI{7942}{\mega\watt} en 2025. En el mismo período, la participación conjunta de la generación solar y eólica pasó de un valor cercano a cero a 9,5\,\pct{}. Estos cambios aumentan la exigencia de transportar energía y aprovechar la capacidad disponible de la red \parencites{minem2025anuario2024}{minem2025indicadoresDiciembre}{coesMaximaDemanda}{iea2025transmissiongrid}.

Los cables enterrados se usan en redes colectoras de plantas solares y eólicas y en tramos donde el espacio, la servidumbre, la seguridad o la integración urbana dificultan construir líneas aéreas \parencite{cigre2023overheadLines}. Su capacidad de transporte depende de la evacuación del calor hacia el lecho de asiento, el suelo de relleno y el suelo nativo. La ampacidad $\Imax$ es la corriente que un cable puede transportar de manera sostenida, bajo condiciones especificadas, sin que la temperatura máxima del conductor $\Tmax$ supere el límite térmico admisible \parencites[cláusulas~4--5]{iec60287}[3--5]{enescu2020}.

El objeto de estudio es el arreglo de cables--entorno térmico, formado por los cables, su disposición y los materiales circundantes. El entorno es heterogéneo cuando su conductividad térmica $\kx$ varía por humedad, estratos o interfaces. Esa variación cambia el campo $T(x,y)$, del cual se obtienen $\Tmax$ y $\Imax$ \parencites[3--4]{enescu2021}[1,6--7]{khumalo2025}[3]{xing2023}.

El problema aparece cuando una propiedad homogénea equivalente sustituye esa distribución espacial. La simplificación puede ocultar zonas de baja disipación o producir estimaciones conservadoras. Los métodos de campo, como FEM, pueden representar la geometría y los materiales con detalle, pero requieren construir, mallar y verificar cada escenario \parencites[41--42,93--99]{cigre2025}[1--3,12]{aldulaimi2024}.

La investigación propone un artefacto basado en una red neuronal informada por física (PINN) que incorpora la ecuación de conducción de calor y sus condiciones de frontera e interfaz. El artefacto aproxima $T(x,y)$, estima $\Tmax$ e $\Imax$ y conserva registros de verificación. La ciencia del diseño organiza la relación entre problema práctico, construcción tecnológica y evidencia de evaluación \parencites[54--56]{peffers2007}[342--350]{gregor2013}.

El Capítulo I delimita el problema, los objetivos y la metodología. El Capítulo II desarrolla los antecedentes y las bases teóricas. El Capítulo III presenta la solución y reserva la evidencia de desarrollo, resultados, discusión e impacto. El Capítulo IV se destina a las conclusiones y recomendaciones que podrán formularse después de evaluar el artefacto.

